\documentclass[11pt,a4paper]{article}
\usepackage[utf8]{inputenc}
\usepackage[vietnamese]{babel}
\usepackage{amsmath,amssymb}
\usepackage{geometry}
\geometry{a4paper, total={180mm,260mm}, left=15mm, right=15mm, top=18mm, bottom=18mm}

\usepackage{tikz}
\usetikzlibrary{shapes.geometric, arrows.meta, positioning, calc}
\tikzset{
  arr/.style={-{Stealth[scale=0.8]}},
  every node/.append style={align=center}
}
\usepackage{tcolorbox}
\tcbuselibrary{skins,breakable}
\usepackage{fontawesome}
\usepackage{fancyhdr}
\usepackage{enumitem}

\pagestyle{fancy}
\fancyhf{}
\renewcommand{\headrulewidth}{0.5pt}
\renewcommand{\footrulewidth}{0.5pt}
\fancyhead[L]{\small\bfseries\color{violet!85!black}\faBook\ SỔ TAY KIẾN THỨC TOÁN 2}
\fancyhead[R]{\small\bfseries\color{violet!85!black}Lớp Toán Thầy Lĩnh -- loptoanthaylinh.vn}
\fancyfoot[C]{\small\color{gray}-- Sổ tay kiến thức: Bài 21. Luyện tập phép cộng có nhớ trong phạm vi 100 --}

\setlength{\parindent}{0pt}
\setlength{\parskip}{0pt}

\begin{document}

% --- BANNER BÀI HỌC ---
\begin{tcolorbox}[
  enhanced,
  colback=violet!8!white,
  colframe=violet!70!black,
  arc=3mm, boxrule=1.2pt,
  left=3mm, right=3mm, top=2mm, bottom=2mm,
  halign=center
]
  {\LARGE\bfseries\color{violet!90!black}\faBookmark\ BÀI 21: LUYỆN TẬP PHÉP CỘNG CÓ NHỚ TRONG PHẠM VI 100}\\[1mm]
  {\itshape\footnotesize\color{black!75} Rèn luyện kĩ năng đặt tính, tính nhẩm nhanh và giải toán lời văn phép cộng có nhớ trên 1 trang A4}
\end{tcolorbox}

\vspace{-2mm}

\noindent
\begin{minipage}[t]{0.49\textwidth}\vspace{0pt}
\begin{tcolorbox}[
  enhanced,
  colback=violet!4!white,
  colframe=violet!75!black,
  arc=2mm, boxrule=0.8pt,
  left=2mm, right=2mm, top=1.8mm, bottom=1.8mm,
  title={\footnotesize\bfseries\color{white}\faKey\ 1. Tổng hợp các dạng phép cộng có nhớ}
]
{\footnotesize
\begin{itemize}\setlength{\itemsep}{1pt}\setlength{\leftskip}{-2.5mm}
  \item \textbf{Dạng 1: Cộng ra số tròn chục}\\
  $26 + 4 = 30$; $35 + 15 = 50$.
  \item \textbf{Dạng 2: Số có 2 chữ số + Số có 1 chữ số}\\
  $48 + 6 = 54$; $75 + 8 = 83$.
  \item \textbf{Dạng 3: Số có 2 chữ số + Số có 2 chữ số}\\
  $38 + 27 = 65$; $49 + 36 = 85$.
  \item \textbf{Khẩu quyết:}\\
  "Cộng đơn vị trước $\to$ Nhớ sang chục $\to$ Cộng chục thêm nhớ".
\end{itemize}
}
\end{tcolorbox}
\end{minipage}
\hfill
\begin{minipage}[t]{0.49\textwidth}\vspace{0pt}
\begin{tcolorbox}[
  enhanced,
  colback=teal!4!white,
  colframe=teal!75!black,
  arc=2mm, boxrule=0.8pt,
  left=2mm, right=2mm, top=1.8mm, bottom=1.8mm,
  title={\footnotesize\bfseries\color{white}\faSliders\ 2. Kĩ năng tính nhanh và so sánh}
]
{\footnotesize
\begin{itemize}\setlength{\itemsep}{1pt}\setlength{\leftskip}{-2.5mm}
  \item \textbf{Tính chất giao hoán và kết hợp:}\\
  Đổi chỗ các số hạng để nhóm tròn chục:\\
  $25 + 18 + 5 = (25 + 5) + 18 = 30 + 18 = 48$.
  \item \textbf{So sánh hai tổng:}\\
  - Nếu tính được ra số: tính hai vế rồi so sánh.\\
  - Nếu có chung một số hạng: số hạng kia lớn hơn thì tổng lớn hơn ($28 + 15 > 28 + 12$).
  \item \textbf{Tìm số hạng chưa biết dạng nhẩm.}
\end{itemize}
}
\end{tcolorbox}
\end{minipage}

\vspace{-1.5mm}
\begin{tcolorbox}[
  enhanced,
  colback=orange!4!white,
  colframe=orange!80!black,
  arc=2mm, boxrule=0.8pt,
  left=2mm, right=2mm, top=1.5mm, bottom=1.5mm,
  title={\footnotesize\bfseries\color{white}\faSitemap\ 3. Sơ đồ tư duy: Bản đồ các dạng cộng có nhớ phạm vi 100}
]
\centering
\begin{tikzpicture}[x=1.3cm, y=0.8cm]
  \node[rectangle, rounded corners=2pt, draw=violet!80!black, fill=violet!20, font=\scriptsize\bfseries] (c) at (0,0) {PHÉP CỘNG\\CÓ NHỚ $\le 100$};
  \node[rectangle, rounded corners=2pt, draw=teal!80!black, fill=teal!15, font=\scriptsize, left=1cm of c] (d1) {CỘNG TRÒN CHỤC\\$26 + 4 = 30$\\$36 + 24 = 60$};
  \node[rectangle, rounded corners=2pt, draw=orange!80!black, fill=orange!15, font=\scriptsize, above right=0.1cm and 0.8cm of c] (d2) {CỘNG VỚI 1 CHỮ SỐ\\$48 + 5 = 53$};
  \node[rectangle, rounded corners=2pt, draw=green!80!black, fill=green!15, font=\scriptsize, below right=0.1cm and 0.8cm of c] (d3) {CỘNG HAI SỐ 2 CHỮ SỐ\\$47 + 25 = 72$};

  \draw[arr, thick, teal!80!black, <-] (d1) -- (c);
  \draw[arr, thick, orange!80!black, ->] (c) -- (d2);
  \draw[arr, thick, green!80!black, ->] (c) -- (d3);
\end{tikzpicture}
\end{tcolorbox}

\vspace{-1.5mm}
\noindent
\begin{minipage}[t]{0.49\textwidth}\vspace{0pt}
\begin{tcolorbox}[
  enhanced,
  colback=green!4!white,
  colframe=green!60!black,
  arc=2mm, boxrule=0.8pt,
  left=2mm, right=2mm, top=1.8mm, bottom=1.8mm,
  title={\footnotesize\bfseries\color{white}\faCheckCircle\ 4. Ví dụ thực hành chuẩn mực}
]
{\footnotesize
\textbf{1. Tính giá trị biểu thức:}\\
$38 + 24 + 6 = 38 + (24 + 6) = 38 + 30 = \mathbf{68}$.
\textbf{2. Bài toán có lời văn:}\\
Vườn nhà bác Ba có $48$ cây cam, bác trồng thêm $25$ cây cam nữa. Hỏi vườn nhà bác Ba có tất cả bao nhiêu cây cam?\\
\textit{Bài giải:}\\
Vườn nhà bác Ba có tất cả số cây cam là:
$$48 + 25 = 73\text{ (cây cam)}$$
Đáp số: $73$ cây cam.
}
\end{tcolorbox}
\end{minipage}
\hfill
\begin{minipage}[t]{0.49\textwidth}\vspace{0pt}
\begin{tcolorbox}[
  enhanced,
  colback=red!4!white,
  colframe=red!75!black,
  arc=2mm, boxrule=0.8pt,
  left=2mm, right=2mm, top=1.8mm, bottom=1.8mm,
  title={\footnotesize\bfseries\color{white}\faExclamationTriangle\ 5. Mẹo tránh bẫy trắc nghiệm}
]
{\footnotesize
\begin{itemize}\setlength{\itemsep}{1pt}\setlength{\leftskip}{-2.5mm}
  \item \textbf{Cộng dãy 3 số:}\\
  Quan sát xem có 2 số nào cộng lại thành số tròn chục thì gộp trước để tính vừa nhanh vừa không nhầm.
  \item Kiểm tra lại bằng cách trừ: Tổng trừ số hạng này phải ra đúng số hạng kia.
\end{itemize}
}
\end{tcolorbox}
\end{minipage}


\end{document}
